% Exercise 2.2
We have
\begin{equation*}
    \begin{split}
        {(1 + r)}^n &= 2 \\
        \Leftrightarrow \ln{{(1 + r)}^n} &= \ln{2} \\
        \Leftrightarrow n\ln{(1 + r)} &= \ln{2} \\
        \Leftrightarrow n &= \frac{\ln{2}}{\ln{(1 + r)}}.
    \end{split}
\end{equation*}
Then, applying the approximation $\ln{(1 + r)} \approx r$ for small $r$, we
have
\begin{equation*}
    \begin{split}
        n &\approx \frac{\ln{2}}{r} \\
        &= \frac{.69}{r}.
    \end{split}
\end{equation*}
Thus, using $i = 100r$, we have
\begin{equation*}
    \begin{split}
        n &\approx \frac{\ln{2}}{r} \\
        &= \frac{.69}{(i / 100)} \\
        &= \frac{69}{i}.
    \end{split}
\end{equation*}

Now, using the better approximation
$\ln{(1 + r)} \approx r - \frac{1}{2}r^2$, we have for $r \approx .08$ we
conclude that
\begin{equation*}
    \begin{split}
        n &\approx \frac{\ln{2}}{r - \frac{1}{2}r^2} \\
        &= \frac{.69}{.08 - \frac{1}{2}{.08}^2}.
    \end{split}
\end{equation*}
Using the following calculation
\begin{lstlisting}
>>> n = 0.69 / (0.08 - (0.5 * (0.08**2)))
>>> n
8.984374999999998
\end{lstlisting}
and, for $72 / i = 72 / 100r \approx 72 / (100 \cdot .08)$, the following
calculation
\begin{lstlisting}
>>> n = 0.69 / (0.08 - (0.5 * (0.08**2)))
>>> n_approx = 72 / (100 * 0.08)
>>> n_approx
9.0
>>> abs(n_approx - n)
0.015625000000001776
\end{lstlisting}
we can see that $n \approx 72 / i$.
